\documentclass[10pt,a4paper]{amsart}
\usepackage[margin=19mm]{geometry}
\usepackage{amsmath,amssymb,mathtools}
\usepackage[scr=rsfs,cal=euler]{mathalfa}
\usepackage{xfrac}
\usepackage[unicode,hidelinks,bookmarks=false]{hyperref}
\newcommand{\ie}{i.e.\ }
\newcommand{\cf}{cf.\ }
\newcommand{\resp}{respectively\ }
\newcommand{\Subsection}{\subsection}
\providecommand{\nobreakdash}{\nobreak\hbox{-}}
\setlength{\emergencystretch}{2em}
\multlinegap0pt
\usepackage{bm}
\usepackage{bbm}
\usepackage{dsfont}
\usepackage{xspace}
\usepackage{cancel}
\usepackage{mathtools}
\usepackage{amsthm}
\makeatletter
\@ifundefined{lem}{\newtheorem{lem}{Lem}}{}
\@ifundefined{thm}{\newtheorem{thm}{Thm}}{}
\makeatother
\newcommand{\RR}{{{\rm I} \kern -.15em {\rm R} }}
\title[Hegselmann-Krause and Cucker-Smale type models with attracti]{Hegselmann-Krause and Cucker-Smale type models with attractive-repulsive interaction}
\author{Elisa Continelli, Cristina Pignotti}
\date{Source: arXiv:2311.13238v1 (CC BY 4.0)}
\begin{document}
\maketitle

\noindent\emph{Source and licence.} Hegselmann-Krause and Cucker-Smale type models with attractive-repulsive interaction. Version arXiv:2311.13238v1 (CC BY 4.0), distributed under the Creative Commons Attribution 4.0 International licence, as recorded in the arXiv licence metadata for this exact version and shown on the abstract page https://arxiv.org/abs/2311.13238v1. Full rights notice: licenses/arxiv-2311.13238-CC-BY-4.0.txt.\par\smallskip

\noindent\textit{Editorial scope.} Counted unit: the Thm whose source label is \texttt{af} in the article (source0.tex lines 749--754, source label \texttt{af}), delivered together with the article's own complete original proof of it (source0.tex lines 978--1128, one \texttt{proof} environment of 151 lines). No mathematics is written, repaired or re-derived: statement and proof are the article's own text line for line. The proof is detached from the statement in the source: the article states the result at lines 749--754 and prints its proof at lines 978--1128, and the proof environment's own header names the statement, so the association is the article's own. Required context retained uncounted: onoff (lines 86--88); incond (lines 90--92); alpha1 (lines 102--106); tn (lines 108--110); tntilde (lines 662--664); sum1tilde (lines 666--668); tildepsi (lines 672--674); badtilde (lines 703--716); onoffcs (lines 719--723); incondcs (lines 725--729); tnbuoni (lines 732--734); distcs (lines 767--770); dist-1cs (lines 787--790); boundgrowthcs (lines 825--830); boundbc (lines 857--859); n-2cs (lines 882--884). Every definition, display and label of those lines is retained unchanged. Omitted from this package and not redistributed: 1--85, 89--89, 93--101, 107--107, 111--661, 665--665, 669--671, 675--702, 717--718, 724--724, 730--731, 735--748, 755--766, 771--786, 791--824, 831--856, 860--881, 885--977, 1129--1315. Nothing inside a retained excerpt is omitted or altered: every retained body is delivered exactly as the article's lines are written, so no omission receipt of any kind accompanies this package. Citations: the retained excerpts carry no citation command at all, so no bibliography entry of the article is transcribed and no reference is redistributed. Reference adaptation: none was needed and none was made; every reference inside the delivered unit resolves inside it, so no unresolved reference or citation is produced. Printed slips kept literal: no printed word, symbol, hypothesis or number of the counted statement or of its proof was altered. Layout and class: the article's own class is \texttt{article} with packages including graphicx, amscd, amsmath, caption, amsfonts, amssymb, amsthm, mathrsfs, multicol, color, babel, fontenc, textcomp, inputenc; the wrapper is the lane's pinned \texttt{amsart} editorial wrapper at 10pt on A4, which restates the theorem-like environments the retained text needs and adds the article's own macro definitions that the retained text needs (\texttt{\textbackslash RR}), with extra wrapper packages (bm, bbm, dsfont, xspace, cancel, mathtools). The delivered numbering is the unit's own. Assets: no retained span contains a figure environment, a graphics-inclusion command, a PDF-page inclusion, a table or a TikZ picture; no image file is copied into this package, the package declares no asset dependency and no image file is redistributed with it. Licence: version arXiv:2311.13238v1 of this article is distributed under the Creative Commons Attribution 4.0 International licence, as recorded in the arXiv licence metadata for this exact version and shown on the abstract page https://arxiv.org/abs/2311.13238v1. A full rights notice is delivered at licenses/arxiv-2311.13238-CC-BY-4.0.txt. Characters: every retained excerpt is pure ASCII, so the delivered engine, XeLaTeX with its default Latin Modern text font, reports no missing character.
	\begin{equation}\label{onoff}
		\frac{d}{dt}x_{i}(t)=\underset{j:j\neq i}{\sum}\alpha(t) a_{ij}(t)(x_{j}(t)-x_{i}(t)),\quad t>0,	\,\, \forall i=1,\dots,N.
	\end{equation}

	\begin{equation}\label{incond}
		x_{i}(0)=x^{0}_{i}\in \mathbb{R}^{d},\quad \forall i=1,\dots,N.
	\end{equation}

	\begin{equation}\label{alpha1}
		\alpha(0)=1,\quad\alpha(t)=\begin{cases}
			1,\quad &t\in \left(t_{2n},t_{2n+1}\right), \,\ \, \ n\in \mathbb{N}_{0},\\-1,\quad &t\in \left[t_{2n+1},t_{2n+2}\right], \,\, n\in \mathbb{N}_{0},
		\end{cases}
	\end{equation}

	\begin{equation}\label{tn}
		t_{2n+2}-t_{2n+1}<\frac{\ln 2}{K},\quad \forall n\in \mathbb{N}_{0}.
	\end{equation}

	Let $\tilde{\psi}:\RR\rightarrow\RR$ be a positive, bounded, continuous function. Assume that the sequence $\{t_{n}\}_{n}$ of definition \eqref{alpha1} satisfies \begin{equation}\label{tntilde}
		t_{2n+2}-t_{2n+1}<\frac{\ln 2}{\tilde K},\quad \forall n\in \mathbb{N}_0.
	\end{equation}

	\begin{equation}\label{sum1tilde}
		\sum_{p=0}^{\infty}\ln\left(\frac{e^{\tilde K(t_{2p+2}-t_{2p+1})}}{2-e^{\tilde K(t_{2p+2}-t_{2p+1})}}\right)<+\infty,
	\end{equation}

where \begin{equation}\label{tildepsi}
	\tilde{\psi}_{0}:=\min_{\lvert y\rvert\leq M^{0}}\tilde\psi(y),
\end{equation}

\begin{thm}
	Let $\psi:\RR\rightarrow\RR$ be a positive, bounded, continuous function. Assume that the sequence $\{t_{n}\}_{n}$ of definition \eqref{alpha1} satisfies \eqref{tntilde}. Assume also that the following conditions hold:
	\begin{equation}\label{badtilde}
		\sum_{p=0}^{+\infty}(t_{2p+2}-t_{2p+1})<+\infty,
	\end{equation}
	\begin{equation}\label{addcond2}
		\sup_{n\in \mathbb{N}}\left(\frac{e^{\tilde K(t_{2n+2}-t_{2n+1})}}{2-e^{\tilde K(t_{2n+2}-t_{2n+1})}}\max\left\{1-e^{-\tilde K(t_{2n+1}-t_{2n})},1-\frac{\tilde{\psi}_{0}}{\tilde K}(1-e^{-\tilde K(t_{2n+1}-t_{2n})})\right\}\right)=c<1,
	\end{equation}
where $\tilde{\psi}_0$ is the positive constant in \eqref{tildepsi}. Then, every solution $\{x_{i}\}_{i=1,\dots,N}$ to \eqref{onoff} with the initial conditions \eqref{incond} satisfies the following exponential decay estimate
	\begin{equation}\label{expestnb}
		d(t)\leq e^{-\gamma\left(t-\frac{\ln2}{\tilde K}-T\right) }d(0),\quad \forall t\geq 0,
	\end{equation}
	for two suitable positive constants $\gamma$, independent of $N$, and $T$.
\end{thm}

\begin{equation}\label{onoffcs}
	\begin{cases}
		\frac{d}{dt}x_{i}(t)=v_{i}(t),&\quad t>0,	\,\,\forall i=1,\dots,N,\\\frac{d}{dt}v_{i}(t)=\underset{j:j\neq i}{\sum}\alpha(t) a_{ij}(t)(x_{j}(t)-x_{i}(t)),&\quad t>0,	\,\,\forall i=1,\dots,N,
	\end{cases}
\end{equation}

\begin{equation}\label{incondcs}
	\begin{cases}
		x_{i}(0)=x^{0}_{i}\in \mathbb{R}^{d},&\quad \forall i=1,\dots,N,\\v_{i}(0)=v^{0}_{i}\in \mathbb{R}^{d},&\quad \forall i=1,\dots,N.
	\end{cases}		
\end{equation}

\begin{equation}\label{tnbuoni}
	t_{2n+1}-t_{2n}>\frac{1}{\tilde K},\quad \forall n\in \mathbb{N}_0.
\end{equation}

	For each $n\in \mathbb{N}_{0}$ and $i,j=1,\dots,N$, we get \begin{equation}
		\label{distcs}
		\lvert v_{i}(s)-v_{j}(t)\rvert\leq d_{V}(t_{2n}), \quad\forall s,t\in [t_{2n},t_{2n+1}].
	\end{equation} 

	For each $n\in \mathbb{N}_{0}$ and $i,j=1,\dots,N$, we get \begin{equation}
		\label{dist-1cs}
		\lvert v_{i}(s)-v_{j}(t)\rvert\leq d_{V}(t_{2n+2}), \quad\forall s,t\in [t_{2n+1},t_{2n+2}].
	\end{equation} 

\begin{lem}
	For all $n\in \mathbb{N}_{0}$, we have that
	\begin{equation}\label{boundgrowthcs}
		d_{V}(t_{2n+2})\leq \frac{e^{\tilde K(t_{2n+2}-t_{2n+1})}}{2-e^{\tilde K(t_{2n+2}-t_{2n+1})}}d_{V}(t_{2n+1}).
	\end{equation}
\end{lem} 

\begin{equation}\label{boundbc}
	t_{n+1}-t_{n}\leq T,\quad \forall n\in \mathbb{N}_0.
\end{equation} 

	\begin{equation}\label{n-2cs}
		d_{V}(t_{2n+1})\leq \left(1-\int_{t_{2n}}^{t_{2n+1}}\phi(s)ds\right)d_{V}(t_{2n}).
	\end{equation}

\begin{thm}\label{af}
	Let $\tilde\psi:\RR^d\times \RR^d\rightarrow\RR$ be a positive, bounded, continuous function satisfying \begin{equation}\label{infint}
		\int_{0}^{\infty}\min_{r\in [0,x]}\tilde \psi(r)dx=+\infty.
	\end{equation} Assume that the sequence $\{t_{n}\}_{n}$ of definition \eqref{alpha1} satisfies \eqref{tntilde}, \eqref{tnbuoni}. Moreover, assume that the condition \eqref{sum1tilde} holds.
	Then, every solution $\{x_{i},v_{i}\}_{i=1,\dots,N}$ to \eqref{onoffcs} with the initial conditions \eqref{incondcs} exhibits asymptotic flocking.
	\end{thm}

\begin{proof}[Proof of Theorem \ref{af}]
	Let $\{(x_{i},v_{i})\}_{i=1,\dots,N}$  be solution to \eqref{onoffcs} under the initial conditions \eqref{incondcs}. We define the function $\mathcal{D}:[0,\infty)\rightarrow [0,\infty),$ 
	$$\mathcal{D}(t):=\begin{cases}
		d_{V}(0), &t=0,\\\left(1-\int_{t_{2n}}^{t}\phi(s)ds\right)d_V(t_{2n}), &t\in (t_{2n},t_{2n+1}],\,n\in \mathbb{N}_0,\\\left(1-\int_{t_{2n+1}}^{t}\phi(s)ds\right)d_V(t_{2n+2}),&t\in (t_{2n+1},t_{2n+2}],\,n\in \mathbb{N}_0.
	\end{cases}$$
	By construction, $\mathcal{D}$ is piecewise continuous. Indeed, $\mathcal{D}$ is continuous everywhere except at points $t_n$, $n\in\mathbb{N}_0$. Moreover, for all $n\in \mathbb{N}$, $n\ge 1,$
	\begin{equation}\label{limt2n}
		\lim_{t\to t^{+}_{2n}}\mathcal{D}(t)=d_V(t_{2n})\geq \left(1-\int_{t_{2n-1}}^{t_{2n}}\phi(s)ds\right)d_V(t_{2n})=\mathcal{D}(t_{2n}),
	\end{equation}
	\begin{equation}\label{limt2n+1}
		\begin{array}{l}
			\displaystyle{\lim_{t\to t^{+}_{2n+1}}\mathcal{D}(t)=d_V(t_{2n+2})\leq \frac{e^{\tilde K(t_{2n+2}-t_{2n+1})}}{2-e^{\tilde K(t_{2n+2}-t_{2n+1})}}d_V(t_{2n+1})}\\
\displaystyle{\hspace{1,8 cm}\leq \frac{e^{\tilde K(t_{2n+2}-t_{2n+1})}}{2-e^{\tilde K(t_{2n+2}-t_{2n+1})}}\left(1-\int_{t_{2n}}^{t_{2n+1}}\phi(s)ds\right)d_V(t_{2n})}\\
\displaystyle{
\hspace{2,5 cm}
=\frac{e^{\tilde K(t_{2n+2}-t_{2n+1})}}{2-e^{\tilde K(t_{2n+2}-t_{2n+1})}}\mathcal{D}(t_{2n+1}).}
		\end{array}
	\end{equation}
	Also, $\mathcal{D}$ is nonincreasing in all intervals of the form $(t_{n},t_{n+1}]$, $n\in \mathbb{N}_0$.
	\\Now, notice that, for almost all times $t,$ 
$$\frac{d}{dt}\max_{s\in [0,t]}d_{X}(s)\leq \left\lvert \frac{d}{dt}d_{X}(t)\right\rvert,$$
	since $\max_{s\in [0,t]}d_{X}(s)$ is constant or increases like $d_{X}(t)$. Moreover, for almost all times $$\left\lvert \frac{d}{dt}d_{X}(t)\right\rvert\leq d_{V}(t).$$
	Therefore, for almost all times
	\begin{equation}\label{maxcs}
		\frac{d}{dt}\max_{s\in [0,t]}d_{X}(s)\leq \left\lvert \frac{d}{dt}d_{X}(t)\right\rvert\leq  d_{V}(t).
	\end{equation}
	Next, we define the function $\mathcal{L}:[0,\infty)\rightarrow [0,\infty)$ as follows: $$\mathcal{L}(t):=\mathcal{D}(t)+\int_{0}^{\underset{s\in [0,t]}{\max}d_{X}(s)}\min\left\{e^{- \tilde KT}\min_{\sigma\in [0,r]}\tilde \psi(\sigma),\frac{e^{-\tilde KT}}{T}\right\}dr,$$
	for all $t\geq0$. By definition, $\mathcal{L}$ is piecewise continuous, i.e. $\mathcal{L}$ is continuous everywhere except at points $t_n$, $n\in \mathbb{N}_0$. \\In addition, for each $n\in \mathbb{N}$ and for all $t\in(t_{2n},t_{2n+1}) $, we have that $$\frac{d}{dt}\mathcal{L}(t)=\frac{d}{dt}\mathcal{D}(t)+\min\left\{e^{-\tilde KT}\tilde \psi_{t},\frac{e^{-\tilde KT}}{T}\right\}\frac{d}{dt}\underset{s\in [0,t]}{\max}d_{X}(s)$$$$=\frac{d}{dt}\mathcal{D}(t)+\phi(t)\frac{d}{dt}\underset{s\in [0,t]}{\max}d_{X}(s),$$
	and from \eqref{maxcs} we get $$\begin{array}{l}
		\vspace{0.3cm}\displaystyle{\frac{d}{dt}\mathcal{L}(t)\leq \frac{d}{dt}\mathcal{D}(t)+\phi(t)d_{V}(t)}\\
		\vspace{0.3cm}\displaystyle{\hspace{1.3cm}=-\phi(t)d_V(t_{2n})+\phi(t)d_{V}(t)}\\
		\displaystyle{\hspace{1.3cm}=\phi(t)(d_V(t)-d_V(t_{2n})).}
	\end{array}$$
	Thus, since $d_V(t)\leq d_V(t_{2n})$ from \eqref{distcs}, we can deduce that
	$$\frac{d}{dt}\mathcal{L}(t)\leq 0,\quad \forall t\in (t_{2n},t_{2n+1}).$$
	As a consequence, for all $t_{2n}<s<t\leq t_{2n+1}$, it comes that
	$$\mathcal{L}(t)\leq \mathcal{L}(s).$$ 
	Letting $s\to t_{2n}^+$, we get
	$$\mathcal{L}(t)\leq \lim_{s\to t_{2n}^+}\mathcal{L}(s)=\lim_{s\to t_{2n}^+}\mathcal{D}(s)+\int_{0}^{\underset{s\in [0,t_{2n}]}{\max}d_{X}(s)}\min\left\{e^{-\tilde KT}\min_{\sigma\in [0,r]}\tilde \psi(\sigma),\frac{e^{-\tilde KT}}{T}\right\}dr,$$
	for all $t\in (t_{2n},t_{2n+1}]$. Thus, using \eqref{tn}, \eqref{distcs} and \eqref{limt2n}, we can write
	$$\begin{array}{l}
		\vspace{0.3cm}\displaystyle{\mathcal{L}(t)\leq d_V(t_{2n})+\int_{0}^{\underset{s\in [0,t_{2n}]}{\max}d_{X}(s)}\min\left\{e^{-\tilde KT}\min_{\sigma\in [0,r]}\tilde\psi(\sigma),\frac{e^{-\tilde KT}}{T}\right\}dr}\\
		\vspace{0.3cm}\displaystyle{\hspace{0.5cm}=\frac{1}{\left(1-\int_{t_{2n-1}}^{t_{2n}}\phi(s)ds\right)}\mathcal{D}(t_{2n})+\int_{0}^{\underset{s\in [0,t_{2n}]}{\max}d_{X}(s)}\min\left\{e^{-\tilde KT}\min_{\sigma\in [0,r]}\tilde\psi(\sigma),\frac{e^{-\tilde KT}}{T}\right\}dr}\\
		\vspace{0.3cm}\displaystyle{\hspace{0.5cm}\leq \frac{1}{\left(1-\int_{t_{2n-1}}^{t_{2n}}\phi(s)ds\right)}\left[\mathcal{D}(t_{2n})+\int_{0}^{\underset{s\in [0,t_{2n}]}{\max}d_{X}(s)}\min\left\{e^{-\tilde KT}\min_{\sigma\in [0,r]}\tilde\psi(\sigma),\frac{e^{-\tilde KT}}{T}\right\}dr\right]}\\
		\displaystyle{\hspace{0.5cm}\leq \frac{1}{\left(1-\int_{t_{2n-1}}^{t_{2n}}\phi(s)ds\right)}\mathcal{L}(t_{2n}),}
	\end{array}$$
	for all $t\in (t_{2n},t_{2n+1}]$. So,
	\begin{equation}\label{lia1}
		\mathcal{L}(t)\leq \frac{1}{\left(1-\int_{t_{2n-1}}^{t_{2n}}\phi(s)ds\right)}\mathcal{L}(t_{2n}),\quad \forall t\in [t_{2n},t_{2n+1}].
	\end{equation}
	On the other hand, for all $t\in (t_{2n+1},t_{2n+2})$, we have that 
	$$\frac{d}{dt}\mathcal{L}(t)=\frac{d}{dt}\mathcal{D}(t)+\min\left\{e^{-\tilde KT}\tilde\psi_{t},\frac{e^{-\tilde KT}}{T}\right\}\frac{d}{dt}\underset{s\in [0,t]}{\max}d_{X}(s)$$$$=\frac{d}{dt}\mathcal{D}(t)+\phi(t)\frac{d}{dt}\underset{s\in [0,t]}{\max}d_{X}(s),$$	
	and from \eqref{maxcs} we get $$\begin{array}{l}
		\vspace{0.3cm}\displaystyle{\frac{d}{dt}\mathcal{L}(t)\leq \frac{d}{dt}\mathcal{D}(t)+\phi(t)d_{V}(t)}\\
		\vspace{0.3cm}\displaystyle{\hspace{1.3cm}=-\phi(t)d_V(t_{2n+2})+\phi(t)d_{V}(t)}\\
		\displaystyle{\hspace{1.3cm}=\phi(t)(d_V(t)-d_V(t_{2n+2})).}
	\end{array}$$
	Thus, since $d_V(t)\leq d_V(t_{2n+2})$ from \eqref{dist-1cs}, we can deduce that
	$$\frac{d}{dt}\mathcal{L}(t)\leq 0,\quad \forall t\in (t_{2n+1},t_{2n+2}).$$
	As a consequence, for all $t_{2n+1}<s<t\leq t_{2n+2}$, it comes that
	$$\mathcal{L}(t)\leq \mathcal{L}(s).$$ 
	Letting $s\to t_{2n+1}^+$, we get
	$$\mathcal{L}(t)\leq \lim_{s\to t_{2n+1}^+}\mathcal{L}(s)=\lim_{s\to t_{2n+1}^+}D(s)+\int_{0}^{\underset{s\in [0,t_{2n+1}]}{\max}d_{X}(s)}\min\left\{e^{-\tilde KT}\min_{\sigma\in [0,r]}\tilde\psi(\sigma),\frac{e^{-\tilde KT}}{T}\right\}dr,$$
	for all $t\in (t_{2n+1},t_{2n+2}]$. Thus, using \eqref{limt2n+1}, we can write
	$$\begin{array}{l}
		\vspace{0.3cm}\displaystyle{\mathcal{L}(t)\leq \frac{e^{\tilde K(t_{2n+2}-t_{2n+1})}}{2-e^{\tilde K(t_{2n+2}-t_{2n+1})}}\mathcal{D}(t_{2n+1})+\int_{0}^{\underset{s\in [0,t_{2n+1}]}{\max}d_{X}(s)}\min\left\{e^{-\tilde KT}\min_{\sigma\in [0,r]}\tilde\psi(\sigma),\frac{e^{-\tilde KT}}{T}\right\}dr}\\	
		\displaystyle{\hspace{4.5cm}\leq \frac{e^{\tilde K(t_{2n+2}-t_{2n+1})}}{2-e^{\tilde K(t_{2n+2}-t_{2n+1})}}\mathcal{L}(t_{2n+1}),}
	\end{array}$$
	for all $t\in (t_{2n+1},t_{2n+2}]$. 
	Therefore, 
	\begin{equation}\label{lia2}
		\mathcal{L}(t)\leq \frac{e^{\tilde K(t_{2n+2}-t_{2n+1})}}{2-e^{\tilde K(t_{2n+2}-t_{2n+1})}}\mathcal{L}(t_{2n+1}),\quad \forall t\in [t_{2n+1},t_{2n+2}].
	\end{equation}
	Now, combining \eqref{lia1} and \eqref{lia2}, it turns out that
	\begin{equation}\label{lia3}
		\mathcal{L}(t_{2n+2})\leq \frac{e^{\tilde K(t_{2n+2}-t_{2n+1})}}{2-e^{\tilde K(t_{2n+2}-t_{2n+1})}}\frac{1}{\left(1-\int_{t_{2n-1}}^{t_{2n}}\phi(s)ds\right)}\mathcal{L}(t_{2n}),\quad \forall n\in \mathbb{N}.
	\end{equation}
	Thus, thanks to an induction argument, from \eqref{lia3} it follows that 
	\begin{equation}\label{lia4}
			\mathcal{L}(t_{2n+2})\leq \prod_{p=1}^{n}\left(\frac{e^{\tilde K(t_{2p+2}-t_{2p+1})}}{2-e^{\tilde K(t_{2p+2}-t_{2p+1})}}\frac{1}{\left(1-\int_{t_{2p-1}}^{t_{2p}}\phi(s)ds\right)}\right)\mathcal{L}(t_2),
	\end{equation}
	for all $n\in \mathbb{N}$. 
	\\Now, let $t\geq t_4$. Then, there exists $n\geq 2$ such that $t\in [t_{2n},t_{2n+2}]$. As a consequence, if $t\in [t_{2n},t_{2n+1}]$, from \eqref{lia1} and \eqref{lia4} with $n-1\geq 1$, we get
	$$\begin{array}{l}
		\vspace{0.3cm}\displaystyle{\mathcal{L}(t)\leq \frac{1}{1-\int_{t_{2n-1}}^{t_{2n}}\phi(s)ds}\mathcal{L}(t_{2n})\leq \frac{e^{\tilde{K}(t_{2n+2}-t_{2n+1})}}{2-e^{\tilde{K}(t_{2n+2}-t_{2n+1})}}\frac{1}{1-\int_{t_{2n-1}}^{t_{2n}}\phi(s)ds}\mathcal{L}(t_{2n})}\\
		\vspace{0.3cm}\displaystyle{\hspace{0.8cm}\leq \frac{e^{\tilde{K}(t_{2n+2}-t_{2n+1})}}{2-e^{\tilde{K}(t_{2n+2}-t_{2n+1})}}\frac{1}{1-\int_{t_{2n-1}}^{t_{2n}}\phi(s)ds}\prod_{p=1}^{n-1}\left(\frac{e^{\tilde K(t_{2p+2}-t_{2p+1})}}{2-e^{\tilde K(t_{2p+2}-t_{2p+1})}}\frac{1}{\left(1-\int_{t_{2p-1}}^{t_{2p}}\phi(s)ds\right)}\right)\mathcal{L}(t_2)}\\
		\displaystyle{\hspace{0.8cm}=\prod_{p=1}^{n}\left(\frac{e^{\tilde K(t_{2p+2}-t_{2p+1})}}{2-e^{\tilde K(t_{2p+2}-t_{2p+1})}}\frac{1}{\left(1-\int_{t_{2p-1}}^{t_{2p}}\phi(s)ds\right)}\right)\mathcal{L}(t_2).}
	\end{array}$$
On the other hand, if $t\in [t_{2n+1},t_{2n+2}]$, from \eqref{lia1}, \eqref{lia2} and \eqref{lia4} it comes that
$$\begin{array}{l}
	\vspace{0.3cm}\displaystyle{\mathcal{L}(t)\leq \frac{e^{\tilde K(t_{2n+2}-t{2n+1})}}{2-e^{\tilde K(t_{2n+2}-t{2n+1})}}\mathcal{L}(t_{2n+1})\leq \frac{e^{\tilde K(t_{2n+2}-t{2n+1})}}{2-e^{\tilde K(t_{2n+2}-t{2n+1})}}\frac{1}{1-\int_{t_{2n-1}}^{t_{2n}}\phi(s)ds}\mathcal{L}(t_{2n})}\\
	\displaystyle{\hspace{1.5cm}\leq \prod_{p=1}^{n}\left(\frac{e^{\tilde K(t_{2p+2}-t_{2p+1})}}{2-e^{\tilde K(t_{2p+2}-t_{2p+1})}}\frac{1}{\left(1-\int_{t_{2p-1}}^{t_{2p}}\phi(s)ds\right)}\right)\mathcal{L}(t_2).}
\end{array}$$
Thus, for all $t\geq t_4$,
\begin{equation}\label{lt}
	\begin{split}
		\mathcal{L}(t)\leq \prod_{p=1}^{n}\left(\frac{e^{\tilde K(t_{2p+2}-t_{2p+1})}}{2-e^{\tilde K(t_{2p+2}-t_{2p+1})}}\frac{1}{\left(1-\int_{t_{2p-1}}^{t_{2p}}\phi(s)ds\right)}\right)\mathcal{L}(t_2)\\=e^{\sum_{p=1}^{n}\left[\ln\left(\frac{e^{\tilde K(t_{2p+2}-t_{2p+1})}}{2-e^{\tilde K(t_{2p+2}-t_{2p+1})}}\right)+\ln\left(\frac{1}{1-\int_{t_{2p-1}}^{t_{2p}}\phi(s)ds}\right)\right]}\mathcal{L}(t_2).
	\end{split}
\end{equation}
Now, from \eqref{sum1tilde}, $\sum_{p=1}^{+\infty}\ln\left(\frac{e^{\tilde K(t_{2n+2}-t_{2p+1})}}{2-e^{\tilde{K}(t_{2p+2}-t_{2p+1})}}\right)<+\infty$. Also, $\sum_{p=1}^{+\infty}\ln\left(\frac{1}{1-\int_{t_{2p-1}}^{t_{2p}}\phi(s)ds}\right)<+\infty$. Indeed, from \eqref{boundbc}, it turns out that $$\int_{t_{2p-1}}^{t_{2p}}\phi(s)ds\leq \frac{e^{-\tilde K T}}{T}(t_{2p}-t_{2p-1}) ,\quad \forall p\geq 1.$$	
	Then, since \eqref{sum1tilde} holds (which implies \eqref{badtilde}), we have that $t_{2p}-t_{2p-1}\to 0$, as $p\to \infty$, and we can write
	$$
\begin{array}{l}
\displaystyle{\ln\left(\frac{1}{1-\int_{t_{2p-1}}^{t_{2p}}\phi(s)ds}\right)\leq \ln\left(\frac{1}{1-\frac{e^{-\tilde K T}}{T}(t_{2p}-t_{2p-1})}\right)}\\
\displaystyle{\hspace{1.5 cm}
=-\ln \left(1-\frac{e^{-\tilde K T}}{T}(t_{2p}-t_{2p-1})\right)
\sim \frac{e^{-\tilde K T}}{T}(t_{2p}-t_{2p-1}).}
\end{array}
$$
	As a consequence, since $\sum_{p=1}^{+\infty}(t_{2p}-t_{2p-1})<+\infty$ from \eqref{badtilde}, it holds $\sum_{p=0}^{+\infty}\ln\left(\frac{1}{1-\int_{t_{2p-1}}^{t_{2p}}\phi(s)ds}\right)<+\infty$.
	\\So, setting $$C:=e^{\sum_{p=0}^{\infty}\left[\ln\left(\frac{e^{\tilde K(t_{2n+2}-t_{2p+1})}}{2-e^{\tilde K(t_{2p+2}-t_{2p+1})}}\right)+\ln\left(\frac{1}{1-\int_{t_{2p-1}}^{t_{2p}}\phi(s)ds}\right)\right]}\mathcal{L}(t_2),$$
	taking into account of \eqref{lt}, we can conclude that
	\begin{equation}\label{lia5}
		\mathcal{L}(t)\leq C,\quad \forall t\geq t_4.
	\end{equation}
	As a consequence, for all $t\geq t_4$, by definition of $\mathcal{L}$,
	$$\int_{0}^{\underset{s\in [0,t]}{\max}d_{X}(s)}\min\left\{e^{-\tilde KT}\min_{\sigma\in [0,r]}\tilde\psi(\sigma),\frac{e^{-\tilde KT}}{T}\right\}dr\leq \mathcal{L}(t)\leq C.$$ 
	Letting $t\to \infty$ in the above inequality, we finally get \begin{equation}\label{lim2}
		\int_{0}^{\underset{s\in [0,\infty)}{\sup}d_{X}(s)}\min\left\{e^{-\tilde KT}\min_{\sigma\in [0,r]}\tilde\psi(\sigma),\frac{e^{-\tilde KT}}{T}\right\}dr\leq C.
	\end{equation} 
	Then, since the function $\tilde \psi$ satisfies  \eqref{infint}, from \eqref{lim2}, there exists a positive constant $d^{*}$ such that \begin{equation}\label{firstcond}
		\underset{s\in [0,\infty)}{\sup}d_{X}(s)\leq d^{*}.
	\end{equation}
	Now, we define
	$$\phi^{*}:=\min\left\{e^{-\tilde KT}\psi_{*},\frac{e^{-\tilde KT}}{T}\right\},$$
	where $$\psi_{*}=\min_{r\in[0,d^{*}]}\tilde\psi(r).$$
	Note that $\phi^{*}>0$, being $\tilde\psi$ a positive function. Also, from \eqref{firstcond}, it comes that $$\psi_{*}\leq \min\left\{\tilde\psi(r):r\in \left[0,\max_{s\in [0,t]}d_{X}(s)\right]\right\}=\tilde\psi_{t},$$
	for all $t\geq 0$. Thus, we get $$\phi^{*}\leq \phi(t), \quad \forall t\geq 0.$$
	As a consequence, for all $n\in \mathbb{N}_0$, $$\int_{t_{2n}}^{t_{2n+1}}\phi(s)ds\geq \phi^*(t_{2n+1}-t_{2n}),$$
	from which 
	$$1-\int_{t_{2n}}^{t_{2n+1}}\phi(s)ds\leq 1-\phi^*(t_{2n+1}-t_{2n}).$$
	So, recalling of inequality \eqref{n-2cs}, we can write
	\begin{equation}\label{n-2phi*}
		d_V(t_{2n+1})\leq \left(1-\phi^*(t_{2n+1}-t_{2n})\right)d_V(t_{2n}),\quad\forall n \in \mathbb{N}_0.
	\end{equation}
	Now, using \eqref{boundgrowthcs} and \eqref{n-2phi*}, we have that
	$$d_V(t_{2n+2})\leq \frac{e^{\tilde K(t_{2n+2}-t_{2n+1})}}{2-e^{\tilde K(t_{2n+2}-t_{2n+1})}}d_V(t_{2n+1})\leq \frac{e^{\tilde K(t_{2n+2}-t_{2n+1})}}{2-e^{\tilde K(t_{2n+2}-t_{2n+1})}}\left(1-\int_{t_{2n}}^{t_{2n+1}}\phi(s)ds\right)d_V(t_{2n})$$
	$$\leq \frac{e^{\tilde K(t_{2n+2}-t_{2n+1})}}{2-e^{\tilde K(t_{2n+2}-t_{2n+1})}}\left(1-\phi^*(t_{2n+1}-t_{2n})\right)d_V(t_{2n}).$$
	Thus, using an induction argument, we get
	$$d_V(t_{2n+2})\leq \prod_{p=0}^{n}\left(\frac{e^{\tilde K(t_{2p+2}-t_{2p+1})}}{2-e^{\tilde K(t_{2p+2}-t_{2p+1})}}\left(1-\phi^*(t_{2p+1}-t_{2p})\right)\right)d_V(0)$$
	$$=e^{\sum_{p=0}^{\infty}\ln \left(\frac{e^{\tilde K(t_{2p+2}-t_{2p+1})}}{2-e^{\tilde K(t_{2p+2}-t_{2p+1})}}\left(1-\phi^*(t_{2p+1}-t_{2p})\right)\right)}d_V(0)$$
	$$=e^{\sum_{p=0}^{\infty}\left[\ln \left(\frac{e^{\tilde K(t_{2p+2}-t_{2p+1})}}{2-e^{\tilde K(t_{2p+2}-t_{2p+1})}}\right)+\ln\left(1-\phi^*(t_{2p+1}-t_{2p})\right)\right]}d_V(0)$$
	Now, $\sum_{p=0}^{\infty}\ln \left(\frac{e^{\tilde K(t_{2p+2}-t_{2p+1})}}{2-e^{\tilde K(t_{2p+2}-t_{2p+1})}}\right)<+\infty$ from \eqref{sum1tilde}. Then, the solution $\{x_{i},v_{i}\}_{i=1,\dots,N}$ exhibits asymptotic flocking if the following condition is satisfied:
	\begin{equation}\label{psi*}
		\sum_{p=0}^{\infty}\ln\left(\left(1-\phi^*(t_{2p+1}-t_{2p})\right)\right)=-\infty.
	\end{equation}
	However, the above condition is guaranteed since, from \eqref{tnbuoni},
	$$1-\phi^*(t_{2p+1}-t_{2p})\leq 1-\frac{\phi^*}{\tilde K}\in (0,1),\quad \forall p\in \mathbb{N}_0.$$
	Thus, $$\sum_{p=0}^{\infty}\ln\left(\left(1-\phi^*(t_{2p+1}-t_{2p})\right)\right)\leq \sum_{p=0}^{\infty}\ln\left(1-\frac{\phi^*}{\tilde K}\right)=-\infty,$$
	from which \eqref{psi*} is fulfilled.
\end{proof}

\end{document}
